\newpage
\section{Wnioski}

Celem projektu było zaimplementowanie listy dwukierunkowej w języku C++ z wykorzystaniem wzorców projektowych Factory oraz Iterator. W ramach projektu utworzono klasy odpowiedzialne za przechowywanie elementów listy, zarządzanie jej zawartością oraz poruszanie się pomiędzy kolejnymi elementami.

Zaimplementowane operacje umożliwiają dodawanie i usuwanie elementów na początku, na końcu oraz pod wskazanym indeksem. Program pozwala również na wyświetlanie zawartości listy w obu kierunkach, przechodzenie do następnego i poprzedniego elementu oraz usuwanie wszystkich elementów.

Podczas realizacji projektu wykorzystano wskaźniki i dynamiczne zarządzanie pamięcią. Istotnym elementem implementacji było prawidłowe aktualizowanie połączeń pomiędzy węzłami, aby zachować strukturę listy po wykonaniu operacji dodawania i usuwania. Zastosowanie destruktora oraz metody \texttt{clear()} umożliwia zwalnianie pamięci zajmowanej przez elementy listy.

Projekt pozwolił na praktyczne zastosowanie programowania obiektowego, podziału kodu na pliki nagłówkowe i źródłowe oraz wykorzystanie wzorców projektowych. W rezultacie powstała struktura danych obsługująca podstawowe operacje na liście dwukierunkowej, stanowiąca przykład zastosowania języka C++ do tworzenia własnych struktur danych.
\begin{table}[h]
\centering
\caption{Podsumowanie funkcjonalności projektu}
\begin{tabular}{|l|l|}
\hline
\textbf{Funkcjonalność} & \textbf{Opis} \\ \hline
Dodawanie elementów & Początek, koniec i wybrany indeks \\ \hline
Usuwanie elementów & Początek, koniec i wybrany indeks \\ \hline
Wyświetlanie listy & Kolejność od początku do końca \\ \hline
Wyświetlanie odwrotne & Kolejność od końca do początku \\ \hline
Iterator & Przejście do następnego i poprzedniego elementu \\ \hline
Czyszczenie listy & Usunięcie wszystkich elementów \\ \hline
NodeFactory & Tworzenie nowych węzłów listy \\ \hline
Zarządzanie pamięcią & Zwalnianie pamięci usuwanych węzłów \\ \hline
\end{tabular}
\label{tab:podsumowanie-funkcjonalnosci}
\end{table}